\section{Introducción}
En el Laboratorio 3 de Electrónica Digital 2 se abordó el diseño de circuitos síncronos mediante VHDL sobre una FPGA Spartan-7, conforme a la consigna de la asignatura \cite{consigna}. Se plantearon dos desarrollos relacionados: en la Parte A, un divisor de frecuencia y un contador binario de diez bits; en la Parte B, una máquina de estados finitos (FSM) encargada de secuenciar operaciones lógicas y aritméticas. La habilitación temporal diseñada para la primera parte se reutilizó conceptualmente en la segunda.

La validación se dividió en simulación funcional e implementación física. \completar{Precisar aquí, en pasado, cuáles etapas se completaron realmente y cuáles quedaron sin ejecutar.}

\subsection{Objetivos}
Se establecieron los siguientes objetivos:
\begin{itemize}
\item Describir procesos síncronos en VHDL mediante \texttt{rising\_edge} y reset.
\item Calcular e implementar un divisor basado en un contador módulo \texttt{MAX}.
\item Implementar un contador ascendente con habilitación \emph{clock enable}.
\item Describir una FSM de Moore mediante estados enumerados y tres procesos.
\item Integrar la FSM como unidad de control y un \emph{datapath} combinacional.
\item Contrastar las formas de onda simuladas con el comportamiento de la FPGA.
\end{itemize}

